% =====================================================================
% Optional supporting material, version V4.8 major revision (2026-08-20).
% Audit-only companion: no result required to support the manuscript's main
% claims appears only in this document. It provides complete cell
% construction, aggregation definitions, per-seed/SNR/class results,
% architecture details, frozen-control analyses, convergence histories,
% prospective-freeze provenance, and artifact paths for results summarized in
% the manuscript.
% =====================================================================
\documentclass[letterpaper,onecolumn,11pt,journal]{IEEEtran}
\usepackage[margin=1.1in]{geometry}
\usepackage[T1]{fontenc}
\usepackage{amsmath,amssymb}
\usepackage{booktabs}
\usepackage{graphicx}
\usepackage[hidelinks]{hyperref}
\usepackage{url}
\urlstyle{tt}
\makeatletter
\g@addto@macro\UrlBreaks{\UrlOrds\do\_\do\/\do\-\do\.}
\makeatother
\setlength{\Urlmuskip}{0mu plus 1mu}
\sloppy

\newcommand{\pp}{percentage points}
% Generated by paper_data_layer/build_paper_numbers.py -- do not edit by hand.
% Every value carries an evidence class; see outputs/paper_numbers.json.

\newcommand{\AbsAFiveClean}{0.5005}% sealed
\newcommand{\AbsAFiveHard}{0.4406}% sealed
\newcommand{\AbsAFiveId}{0.5654}% sealed
\newcommand{\AbsAZeroClean}{0.4705}% sealed
\newcommand{\AbsAZeroHard}{0.3949}% sealed
\newcommand{\AbsAZeroId}{0.5102}% sealed
\newcommand{\AbsCSSLClean}{0.5974}% sealed
\newcommand{\AbsCSSLHard}{0.3915}% sealed
\newcommand{\AbsCSSLId}{0.5158}% sealed
\newcommand{\AbsIQFormerClean}{0.6919}% sealed
\newcommand{\AbsIQFormerHard}{0.4872}% sealed
\newcommand{\AbsIQFormerId}{0.6880}% sealed
\newcommand{\AbsMCLDNNClean}{0.6414}% sealed
\newcommand{\AbsMCLDNNHard}{0.4599}% sealed
\newcommand{\AbsMCLDNNId}{0.6439}% sealed
\newcommand{\BreakEvenIqformerSirFifteen}{0.11}% exploratory
\newcommand{\BreakEvenIqformerSirTen}{0.47}% exploratory
\newcommand{\BreakEvenMcldnnSirFifteen}{0.06}% exploratory
\newcommand{\BreakEvenMcldnnSirTen}{0.29}% exploratory
\newcommand{\CapacityAFiveDoublingsVsAZero}{2.22}% sealed
\newcommand{\CapacityImpliedShareHigh}{3.2}% exploratory
\newcommand{\CapacityImpliedShareHighPct}{71}% exploratory
\newcommand{\CapacityImpliedShareLow}{2.8}% exploratory
\newcommand{\CapacityImpliedShareLowPct}{60}% exploratory
\newcommand{\CapacityLCleanVsIqformerCandidate}{51.16}% exploratory
\newcommand{\CapacityLCleanVsIqformerDiff}{-18.14}% exploratory
\newcommand{\CapacityLCleanVsIqformerHigh}{-17.00}% exploratory
\newcommand{\CapacityLCleanVsIqformerLow}{-19.34}% exploratory
\newcommand{\CapacityLCleanVsIqformerMagnitude}{18.14}% exploratory
\newcommand{\CapacityLCleanVsIqformerReference}{69.30}% exploratory
\newcommand{\CapacityLCleanVsIqformerSeeds}{5}% exploratory
\newcommand{\CapacityLCleanVsSCandidate}{51.16}% exploratory
\newcommand{\CapacityLCleanVsSDiff}{1.20}% exploratory
\newcommand{\CapacityLCleanVsSHigh}{1.82}% exploratory
\newcommand{\CapacityLCleanVsSLow}{0.63}% exploratory
\newcommand{\CapacityLCleanVsSReference}{49.96}% exploratory
\newcommand{\CapacityLCleanVsSSeeds}{5}% exploratory
\newcommand{\CapacityLHardVsIqformerCandidate}{46.78}% exploratory
\newcommand{\CapacityLHardVsIqformerDiff}{-2.84}% exploratory
\newcommand{\CapacityLHardVsIqformerHigh}{-1.31}% exploratory
\newcommand{\CapacityLHardVsIqformerLow}{-4.30}% exploratory
\newcommand{\CapacityLHardVsIqformerMagnitude}{2.84}% exploratory
\newcommand{\CapacityLHardVsIqformerReference}{49.61}% exploratory
\newcommand{\CapacityLHardVsIqformerSeeds}{5}% exploratory
\newcommand{\CapacityLHardVsMCandidate}{46.78}% exploratory
\newcommand{\CapacityLHardVsMDiff}{1.24}% exploratory
\newcommand{\CapacityLHardVsMHigh}{2.37}% exploratory
\newcommand{\CapacityLHardVsMLow}{0.31}% exploratory
\newcommand{\CapacityLHardVsMReference}{45.54}% exploratory
\newcommand{\CapacityLHardVsMSeeds}{5}% exploratory
\newcommand{\CapacityLHardVsMcldnnCandidate}{46.78}% exploratory
\newcommand{\CapacityLHardVsMcldnnDiff}{1.37}% exploratory
\newcommand{\CapacityLHardVsMcldnnHigh}{2.72}% exploratory
\newcommand{\CapacityLHardVsMcldnnLow}{0.03}% exploratory
\newcommand{\CapacityLHardVsMcldnnReference}{45.41}% exploratory
\newcommand{\CapacityLHardVsMcldnnSeeds}{5}% exploratory
\newcommand{\CapacityLHardVsSCandidate}{46.78}% exploratory
\newcommand{\CapacityLHardVsSDiff}{3.18}% exploratory
\newcommand{\CapacityLHardVsSHigh}{3.96}% exploratory
\newcommand{\CapacityLHardVsSLow}{2.44}% exploratory
\newcommand{\CapacityLHardVsSReference}{43.59}% exploratory
\newcommand{\CapacityLHardVsSSeeds}{5}% exploratory
\newcommand{\CapacityLMacs}{249.1}% audit
\newcommand{\CapacityLParameters}{233244}% exploratory
\newcommand{\CapacityLSevereVsIqformerCandidate}{40.30}% exploratory
\newcommand{\CapacityLSevereVsIqformerDiff}{14.07}% exploratory
\newcommand{\CapacityLSevereVsIqformerHigh}{18.51}% exploratory
\newcommand{\CapacityLSevereVsIqformerLow}{9.55}% exploratory
\newcommand{\CapacityLSevereVsIqformerReference}{26.23}% exploratory
\newcommand{\CapacityLSevereVsIqformerSeeds}{5}% exploratory
\newcommand{\CapacityMCleanVsSCandidate}{51.06}% exploratory
\newcommand{\CapacityMCleanVsSDiff}{1.10}% exploratory
\newcommand{\CapacityMCleanVsSHigh}{1.74}% exploratory
\newcommand{\CapacityMCleanVsSLow}{0.44}% exploratory
\newcommand{\CapacityMCleanVsSReference}{49.96}% exploratory
\newcommand{\CapacityMCleanVsSSeeds}{5}% exploratory
\newcommand{\CapacityMHardVsSCandidate}{45.54}% exploratory
\newcommand{\CapacityMHardVsSDiff}{1.94}% exploratory
\newcommand{\CapacityMHardVsSHigh}{2.80}% exploratory
\newcommand{\CapacityMHardVsSLow}{1.07}% exploratory
\newcommand{\CapacityMHardVsSReference}{43.59}% exploratory
\newcommand{\CapacityMHardVsSSeeds}{5}% exploratory
\newcommand{\CapacityMMacs}{104.1}% audit
\newcommand{\CapacityMParameters}{99596}% exploratory
\newcommand{\CapacityRatioAFiveToAZero}{4.67}% sealed
\newcommand{\CapacityRatioLargeToShort}{5.90}% exploratory
\newcommand{\CapacityRatioMidToShort}{2.52}% exploratory
\newcommand{\CapacitySlopeLargePerDoubling}{1.24}% exploratory
\newcommand{\CapacitySlopeMidPerDoubling}{1.46}% exploratory
\newcommand{\ComponentEffectBound}{1.36}% exploratory
\newcommand{\ConfirmatoryFamilyGatePassed}{false}% sealed
\newcommand{\CoverageCleanCandidate}{50.20}% exploratory
\newcommand{\CoverageCleanDiff}{0.15}% exploratory
\newcommand{\CoverageCleanHigh}{0.63}% exploratory
\newcommand{\CoverageCleanLow}{-0.31}% exploratory
\newcommand{\CoverageCleanReference}{50.05}% exploratory
\newcommand{\CoverageCleanSeeds}{10}% exploratory
\newcommand{\CoverageHardCandidate}{43.76}% exploratory
\newcommand{\CoverageHardDiff}{-0.30}% exploratory
\newcommand{\CoverageHardHigh}{0.27}% exploratory
\newcommand{\CoverageHardLow}{-0.91}% exploratory
\newcommand{\CoverageHardReference}{44.06}% exploratory
\newcommand{\CoverageHardSeeds}{10}% exploratory
\newcommand{\DesignResolution}{1.31}% exploratory
\newcommand{\EnvelopeCellCount}{113}% exploratory
\newcommand{\EnvelopeHoldoutCells}{81}% exploratory
\newcommand{\EnvelopeIqformerHoldoutR}{0.694}% exploratory
\newcommand{\EnvelopeIqformerHoldoutRmse}{6.20}% exploratory
\newcommand{\EnvelopeIqformerHoldoutSign}{0.975}% exploratory
\newcommand{\EnvelopeIqformerIntercept}{-16.60}% exploratory
\newcommand{\EnvelopeIqformerOccupancySlope}{11.39}% exploratory
\newcommand{\EnvelopeIqformerRSquared}{0.580}% exploratory
\newcommand{\EnvelopeIqformerSirSlope}{-0.666}% exploratory
\newcommand{\EnvelopeMcldnnHoldoutR}{0.511}% exploratory
\newcommand{\EnvelopeMcldnnHoldoutRmse}{8.05}% exploratory
\newcommand{\EnvelopeMcldnnHoldoutSign}{0.951}% exploratory
\newcommand{\EnvelopeMcldnnIntercept}{-14.28}% exploratory
\newcommand{\EnvelopeMcldnnOccupancySlope}{13.93}% exploratory
\newcommand{\EnvelopeMcldnnRSquared}{0.449}% exploratory
\newcommand{\EnvelopeMcldnnSirSlope}{-0.504}% exploratory
\newcommand{\FusionIqformerControl}{1.49}% exploratory
\newcommand{\FusionIqformerExcess}{3.22}% exploratory
\newcommand{\FusionIqformerGain}{4.71}% exploratory
\newcommand{\FusionMcldnnControl}{1.69}% exploratory
\newcommand{\FusionMcldnnExcess}{4.02}% exploratory
\newcommand{\FusionMcldnnGain}{5.70}% exploratory
\newcommand{\HeadlineVsCsslClean}{-9.69}% sealed
\newcommand{\HeadlineVsCsslCleanHigh}{-8.83}% sealed
\newcommand{\HeadlineVsCsslCleanLow}{-10.56}% sealed
\newcommand{\HeadlineVsCsslCombinedOod}{-4.12}% sealed
\newcommand{\HeadlineVsCsslCombinedOodHigh}{-2.96}% sealed
\newcommand{\HeadlineVsCsslCombinedOodLow}{-5.29}% sealed
\newcommand{\HeadlineVsCsslHard}{4.91}% sealed
\newcommand{\HeadlineVsCsslHardHigh}{6.07}% sealed
\newcommand{\HeadlineVsCsslHardLow}{3.78}% sealed
\newcommand{\HeadlineVsCsslHeldoutChannel}{1.13}% sealed
\newcommand{\HeadlineVsCsslHeldoutChannelHigh}{2.30}% sealed
\newcommand{\HeadlineVsCsslHeldoutChannelLow}{0.01}% sealed
\newcommand{\HeadlineVsCsslId}{4.96}% sealed
\newcommand{\HeadlineVsCsslIdHigh}{6.13}% sealed
\newcommand{\HeadlineVsCsslIdLow}{3.78}% sealed
\newcommand{\HeadlineVsCsslUnseenJammer}{-4.14}% sealed
\newcommand{\HeadlineVsCsslUnseenJammerHigh}{-3.08}% sealed
\newcommand{\HeadlineVsCsslUnseenJammerLow}{-5.20}% sealed
\newcommand{\HeadlineVsCsslUnseenSpeed}{1.54}% sealed
\newcommand{\HeadlineVsCsslUnseenSpeedHigh}{2.64}% sealed
\newcommand{\HeadlineVsCsslUnseenSpeedLow}{0.45}% sealed
\newcommand{\LatencyAFive}{3.172}% sealed
\newcommand{\LatencyAZero}{0.729}% sealed
\newcommand{\LatencyCSSL}{0.786}% sealed
\newcommand{\LatencyIQFormer}{4.180}% sealed
\newcommand{\LatencyMCLDNN}{3.789}% sealed
\newcommand{\MacsAFive}{41.8}% sealed
\newcommand{\MacsAZero}{6.3}% sealed
\newcommand{\MacsCSSL}{155.3}% sealed
\newcommand{\MacsIQFormer}{355.6}% sealed
\newcommand{\MacsMCLDNN}{398.2}% sealed
\newcommand{\MethodEffectDiff}{4.57}% sealed
\newcommand{\MethodEffectSimHigh}{5.49}% sealed
\newcommand{\MethodEffectSimLow}{3.65}% sealed
\newcommand{\ParamsAFive}{39500}% sealed
\newcommand{\ParamsAZero}{8458}% sealed
\newcommand{\ParamsCSSL}{8631948}% sealed
\newcommand{\ParamsIQFormer}{354984}% sealed
\newcommand{\ParamsMCLDNN}{406070}% sealed
\newcommand{\PresenceGateCleanMean}{99.994}% exploratory
\newcommand{\PresenceGateCleanVsAFiveCandidate}{49.75}% exploratory
\newcommand{\PresenceGateCleanVsAFiveDiff}{-0.30}% exploratory
\newcommand{\PresenceGateCleanVsAFiveHigh}{0.08}% exploratory
\newcommand{\PresenceGateCleanVsAFiveLow}{-0.66}% exploratory
\newcommand{\PresenceGateCleanVsAFiveReference}{50.05}% exploratory
\newcommand{\PresenceGateCleanVsAFiveSeeds}{10}% exploratory
\newcommand{\PresenceGateHardMean}{99.991}% exploratory
\newcommand{\PresenceGateHardVsAFiveCandidate}{43.87}% exploratory
\newcommand{\PresenceGateHardVsAFiveDiff}{-0.18}% exploratory
\newcommand{\PresenceGateHardVsAFiveHigh}{0.35}% exploratory
\newcommand{\PresenceGateHardVsAFiveLow}{-0.72}% exploratory
\newcommand{\PresenceGateHardVsAFiveReference}{44.06}% exploratory
\newcommand{\PresenceGateHardVsAFiveSeeds}{10}% exploratory
\newcommand{\ProbeAFiveBelowThreshold}{5}% exploratory
\newcommand{\ProbeAFiveCells}{6}% exploratory
\newcommand{\ProbeStrongAboveThreshold}{18}% exploratory
\newcommand{\ProbeStrongCells}{18}% exploratory
\newcommand{\ProbeThreshold}{0.25}% exploratory
\newcommand{\ProtocolBootstrapDraws}{10000}% sealed
\newcommand{\ProtocolFitCount}{120}% sealed
\newcommand{\ProtocolModelCount}{12}% sealed
\newcommand{\ProtocolSeedCount}{10}% sealed
\newcommand{\ProtocolSplitCount}{11}% sealed
\newcommand{\ProtocolTrainSources}{100000}% sealed
\newcommand{\RetrainedCampaignAFiveLevel}{43.59}% exploratory
\newcommand{\RetrainedCampaignFullLevel}{48.35}% exploratory
\newcommand{\RetrainedCampaignIqzeroLevel}{43.38}% exploratory
\newcommand{\RetrainedCampaignIqzeroVsAFiveDiff}{-0.21}% exploratory
\newcommand{\RetrainedFullPooledLevel}{26.11}% exploratory
\newcommand{\RetrainedFullVsAFiveCiHigh}{7.66}% exploratory
\newcommand{\RetrainedFullVsAFiveCiLow}{3.54}% exploratory
\newcommand{\RetrainedFullVsAFiveDiff}{5.39}% exploratory
\newcommand{\RetrainedIqzeroPooledLevel}{20.18}% exploratory
\newcommand{\RetrainedIqzeroVsAFiveCiHigh}{0.26}% exploratory
\newcommand{\RetrainedIqzeroVsAFiveCiLow}{-1.34}% exploratory
\newcommand{\RetrainedIqzeroVsAFiveDiff}{-0.54}% exploratory
\newcommand{\RetrainedIqzeroVsFullCiHigh}{-4.69}% exploratory
\newcommand{\RetrainedIqzeroVsFullCiLow}{-7.53}% exploratory
\newcommand{\RetrainedIqzeroVsFullDiff}{-5.93}% exploratory
\newcommand{\SelectiveAFiveAtThirty}{94.4}% exploratory
\newcommand{\SelectiveIQFormerAtThirty}{94.9}% exploratory
\newcommand{\SevereCellsInspected}{64}% exploratory
\newcommand{\SevereCellsSurvivingHolm}{12}% exploratory
\newcommand{\SevereControlDecisionPattern}{mixed_representation_and_fusion}% exploratory
\newcommand{\SevereControlDiffCapL}{0.17}% exploratory
\newcommand{\SevereControlDiffCapLHigh}{2.51}% exploratory
\newcommand{\SevereControlDiffCapLLow}{-2.47}% exploratory
\newcommand{\SevereControlDiffCapM}{1.05}% exploratory
\newcommand{\SevereControlDiffCapMHigh}{3.16}% exploratory
\newcommand{\SevereControlDiffCapMLow}{-1.47}% exploratory
\newcommand{\SevereControlDiffIqShuffle}{-4.47}% exploratory
\newcommand{\SevereControlDiffIqShuffleHigh}{-3.36}% exploratory
\newcommand{\SevereControlDiffIqShuffleLow}{-5.58}% exploratory
\newcommand{\SevereControlDiffIqZero}{-7.91}% exploratory
\newcommand{\SevereControlDiffIqZeroHigh}{-5.91}% exploratory
\newcommand{\SevereControlDiffIqZeroLow}{-9.93}% exploratory
\newcommand{\SevereControlDiffNoSideHead}{-13.62}% exploratory
\newcommand{\SevereControlDiffNoSideHeadHigh}{-11.76}% exploratory
\newcommand{\SevereControlDiffNoSideHeadLow}{-15.50}% exploratory
\newcommand{\SevereControlIqShuffleLoss}{8.84}% exploratory
\newcommand{\SevereControlIqShuffleRepairRatio}{2.02}% exploratory
\newcommand{\SevereControlIqZeroLoss}{12.28}% exploratory
\newcommand{\SevereControlIqZeroRepairRatio}{2.81}% exploratory
\newcommand{\SevereControlLevelCapL}{20.89}% exploratory
\newcommand{\SevereControlLevelCapM}{21.77}% exploratory
\newcommand{\SevereControlLevelIqShuffle}{16.25}% exploratory
\newcommand{\SevereControlLevelIqZero}{12.80}% exploratory
\newcommand{\SevereControlLevelNoSideHead}{7.09}% exploratory
\newcommand{\SevereControlMatchedLevelAFive}{20.72}% exploratory
\newcommand{\SevereControlMatchedLevelSidecar}{26.11}% exploratory
\newcommand{\SevereControlNoSideHeadLoss}{17.99}% exploratory
\newcommand{\SevereControlNoSideHeadRepairRatio}{4.12}% exploratory
\newcommand{\SevereGainIqformer}{10.46}% exploratory
\newcommand{\SevereGainIqformerHigh}{14.29}% exploratory
\newcommand{\SevereGainIqformerLow}{6.78}% exploratory
\newcommand{\SevereGainIqformerPositiveSeeds}{10}% exploratory
\newcommand{\SevereGainIqformerSignFlipP}{0.0018}% exploratory
\newcommand{\SevereGainMcldnn}{9.74}% exploratory
\newcommand{\SevereGainMcldnnHigh}{12.14}% exploratory
\newcommand{\SevereGainMcldnnLow}{7.52}% exploratory
\newcommand{\SevereGainMcldnnPositiveSeeds}{10}% exploratory
\newcommand{\SevereGainMcldnnSignFlipP}{0.0018}% exploratory
\newcommand{\SevereHeldFamilyOfdmAFive}{4.60}% exploratory
\newcommand{\SevereHeldFamilyOfdmIqformer}{4.56}% exploratory
\newcommand{\SevereHeldFamilyOfdmMcldnn}{4.92}% exploratory
\newcommand{\SevereHeldFamilyOfdmSidecar}{5.74}% exploratory
\newcommand{\SevereHeldFamilyPulseAFive}{29.03}% exploratory
\newcommand{\SevereHeldFamilyPulseIqformer}{37.45}% exploratory
\newcommand{\SevereHeldFamilyPulseMcldnn}{55.34}% exploratory
\newcommand{\SevereHeldFamilyPulseSidecar}{38.22}% exploratory
\newcommand{\SevereHeldRegimeOneAFive}{21.13}% exploratory
\newcommand{\SevereHeldRegimeOneIqformer}{23.82}% exploratory
\newcommand{\SevereHeldRegimeOneMcldnn}{33.98}% exploratory
\newcommand{\SevereHeldRegimeOneSidecar}{25.34}% exploratory
\newcommand{\SevereHeldRegimeTwoAFive}{20.31}% exploratory
\newcommand{\SevereHeldRegimeTwoIqformer}{22.77}% exploratory
\newcommand{\SevereHeldRegimeTwoMcldnn}{33.52}% exploratory
\newcommand{\SevereHeldRegimeTwoSidecar}{24.83}% exploratory
\newcommand{\SevereHeldSourcesPerRegime}{6400}% exploratory
\newcommand{\SevereHeldSourcesTotal}{12800}% exploratory
\newcommand{\SevereReferenceLeadIqformer}{2.58}% exploratory
\newcommand{\SevereReferenceLeadMcldnn}{13.03}% exploratory
\newcommand{\SevereRepairCiHigh}{5.90}% exploratory
\newcommand{\SevereRepairCiLow}{2.96}% exploratory
\newcommand{\SevereRepairDiff}{4.37}% exploratory
\newcommand{\SevereRepairGapRecoveryIqformer}{1.70}% exploratory
\newcommand{\SevereRepairGapRecoveryMcldnn}{0.34}% exploratory
\newcommand{\SevereRepairLevelAFive}{20.72}% exploratory
\newcommand{\SevereRepairLevelIqformer}{23.29}% exploratory
\newcommand{\SevereRepairLevelMcldnn}{33.75}% exploratory
\newcommand{\SevereRepairLevelSidecar}{25.09}% exploratory
\newcommand{\SevereRepairPositiveSeeds}{10}% exploratory
\newcommand{\SevereRepairPositiveSnrStrata}{8}% exploratory
\newcommand{\SevereRepairRegimeOneDiff}{4.22}% exploratory
\newcommand{\SevereRepairRegimeTwoDiff}{4.52}% exploratory
\newcommand{\SevereRepairResidualGapIqformer}{1.79}% exploratory
\newcommand{\SevereRepairResidualGapMcldnn}{-8.66}% exploratory
\newcommand{\SevereSeedMeanGapIqformer}{-2.58}% exploratory
\newcommand{\SevereSeedMeanGapMcldnn}{-13.03}% exploratory
\newcommand{\SevereTransferGainIqformer}{-3.12}% exploratory
\newcommand{\SevereTransferGainMcldnn}{-12.45}% exploratory
\newcommand{\SevereTransferLeadIqformer}{3.12}% exploratory
\newcommand{\SevereTransferLeadMcldnn}{12.45}% exploratory
\newcommand{\SevereTransferRmseIqformerConstant}{8.11}% exploratory
\newcommand{\SevereTransferRmseIqformerFull}{9.75}% exploratory
\newcommand{\SevereTransferRmseIqformerSirOnly}{9.05}% exploratory
\newcommand{\SevereTransferRmseMcldnnConstant}{15.41}% exploratory
\newcommand{\SevereTransferRmseMcldnnFull}{22.56}% exploratory
\newcommand{\SevereTransferRmseMcldnnSirOnly}{22.29}% exploratory
\newcommand{\SevereTransferSelectorRecall}{0.581}% exploratory
\newcommand{\SevereTransferSignAccIqformer}{0.578}% exploratory
\newcommand{\SevereTransferSignAccMcldnn}{0.406}% exploratory
\newcommand{\SidecarCleanVsAFiveCandidate}{61.82}% exploratory
\newcommand{\SidecarCleanVsAFiveDiff}{11.77}% exploratory
\newcommand{\SidecarCleanVsAFiveHigh}{12.49}% exploratory
\newcommand{\SidecarCleanVsAFiveLow}{11.06}% exploratory
\newcommand{\SidecarCleanVsAFiveReference}{50.05}% exploratory
\newcommand{\SidecarCleanVsAFiveSeeds}{10}% exploratory
\newcommand{\SidecarHardVsAFiveCandidate}{48.43}% exploratory
\newcommand{\SidecarHardVsAFiveDiff}{4.37}% exploratory
\newcommand{\SidecarHardVsAFiveHigh}{5.31}% exploratory
\newcommand{\SidecarHardVsAFiveLow}{3.44}% exploratory
\newcommand{\SidecarHardVsAFiveReference}{44.06}% exploratory
\newcommand{\SidecarHardVsAFiveSeeds}{10}% exploratory
\newcommand{\SidecarHardVsIqformerCandidate}{48.43}% exploratory
\newcommand{\SidecarHardVsIqformerDiff}{-0.29}% exploratory
\newcommand{\SidecarHardVsIqformerHigh}{1.46}% exploratory
\newcommand{\SidecarHardVsIqformerLow}{-1.95}% exploratory
\newcommand{\SidecarHardVsIqformerReference}{48.72}% exploratory
\newcommand{\SidecarHardVsIqformerSeeds}{10}% exploratory
\newcommand{\SidecarHardVsMcldnnCandidate}{48.43}% exploratory
\newcommand{\SidecarHardVsMcldnnDiff}{2.44}% exploratory
\newcommand{\SidecarHardVsMcldnnHigh}{3.49}% exploratory
\newcommand{\SidecarHardVsMcldnnLow}{1.36}% exploratory
\newcommand{\SidecarHardVsMcldnnReference}{45.99}% exploratory
\newcommand{\SidecarHardVsMcldnnSeeds}{10}% exploratory
\newcommand{\SidecarParameters}{46794}% exploratory
\newcommand{\SidecarSevereVsIqformerCandidate}{37.84}% exploratory
\newcommand{\SidecarSevereVsIqformerDiff}{12.09}% exploratory
\newcommand{\SidecarSevereVsIqformerHigh}{15.88}% exploratory
\newcommand{\SidecarSevereVsIqformerLow}{8.32}% exploratory
\newcommand{\SidecarSevereVsIqformerReference}{25.76}% exploratory
\newcommand{\SidecarSevereVsIqformerSeeds}{10}% exploratory
\newcommand{\SidecarSevereVsMcldnnCandidate}{37.84}% exploratory
\newcommand{\SidecarSevereVsMcldnnDiff}{11.37}% exploratory
\newcommand{\SidecarSevereVsMcldnnHigh}{13.70}% exploratory
\newcommand{\SidecarSevereVsMcldnnLow}{9.05}% exploratory
\newcommand{\SidecarSevereVsMcldnnReference}{26.47}% exploratory
\newcommand{\SidecarSevereVsMcldnnSeeds}{10}% exploratory
\newcommand{\TeacherFormDiff}{0.16}% sealed
\newcommand{\TeacherFormSimHigh}{1.08}% sealed
\newcommand{\TeacherFormSimLow}{-0.76}% sealed
\newcommand{\TrainClassesWithCleanWindows}{4}% exploratory
\newcommand{\TrainClassesWithoutCleanWindows}{6}% exploratory
\newcommand{\TransferFailShareCompact}{100}% exploratory
\newcommand{\TransferFailShareHighCapacity}{11}% exploratory
\newcommand{\TriRouteDiff}{0.44}% sealed
\newcommand{\TriRouteSimHigh}{1.36}% sealed
\newcommand{\TriRouteSimLow}{-0.47}% sealed

% Audit-verified envelope constants shared with the manuscript. The v4.1
% audit file renews names that the manuscript defines elsewhere; provide
% placeholders so the renew is safe in this standalone document.
\providecommand{\EnvelopeHoldoutNegShareIQ}{}
\providecommand{\EnvelopeHoldoutNegShareMC}{}
\providecommand{\EnvelopeConstantRmseIQ}{}
\providecommand{\EnvelopeConstantRmseMC}{}
\providecommand{\EnvelopeSkillRTwoIQ}{}
\providecommand{\EnvelopeSkillRTwoMC}{}
\providecommand{\EnvelopeOracleConstRmseIQ}{}
\providecommand{\EnvelopeOracleConstRmseMC}{}
\providecommand{\EnvelopeOracleSkillIQ}{}
\providecommand{\EnvelopeOracleSkillMC}{}
\providecommand{\EnvelopeNoCleanRmseIQ}{}
\providecommand{\EnvelopeNoCleanRmseMC}{}
\providecommand{\EnvelopeNoCleanSkillIQ}{}
\providecommand{\EnvelopeNoCleanSkillMC}{}
\providecommand{\EnvelopeThreeVarRmseIQ}{}
\providecommand{\EnvelopeThreeVarRmseMC}{}
\providecommand{\EnvelopeThreeVarSkillIQ}{}
\providecommand{\EnvelopeThreeVarSkillMC}{}
% V4.1: 80 finite-SIR interfered held cells only. Derived by
% analysis_zero_compute/a14_envelope_v41.py from sealed predictions.
\renewcommand{\EnvelopeHoldoutCells}{80}
\renewcommand{\EnvelopeIqformerHoldoutRmse}{6.23}
\renewcommand{\EnvelopeMcldnnHoldoutRmse}{8.10}
\renewcommand{\EnvelopeIqformerHoldoutSign}{0.975}
\renewcommand{\EnvelopeMcldnnHoldoutSign}{0.950}
\renewcommand{\EnvelopeIqformerIntercept}{-17.08}
\renewcommand{\EnvelopeIqformerOccupancySlope}{14.38}
\renewcommand{\EnvelopeIqformerSirSlope}{-1.036}
\renewcommand{\EnvelopeMcldnnIntercept}{-14.55}
\renewcommand{\EnvelopeMcldnnOccupancySlope}{19.16}
\renewcommand{\EnvelopeMcldnnSirSlope}{-0.894}
\renewcommand{\EnvelopeHoldoutNegShareIQ}{0.963}
\renewcommand{\EnvelopeHoldoutNegShareMC}{0.900}
\renewcommand{\EnvelopeConstantRmseIQ}{8.37}
\renewcommand{\EnvelopeConstantRmseMC}{9.63}
\renewcommand{\EnvelopeSkillRTwoIQ}{0.45}
\renewcommand{\EnvelopeSkillRTwoMC}{0.29}
\renewcommand{\EnvelopeOracleConstRmseIQ}{6.04}
\renewcommand{\EnvelopeOracleConstRmseMC}{6.81}
\renewcommand{\EnvelopeOracleSkillIQ}{-0.06}
\renewcommand{\EnvelopeOracleSkillMC}{-0.42}
\renewcommand{\EnvelopeNoCleanRmseIQ}{6.23}
\renewcommand{\EnvelopeNoCleanRmseMC}{8.10}
\renewcommand{\EnvelopeNoCleanSkillIQ}{0.45}
\renewcommand{\EnvelopeNoCleanSkillMC}{0.29}
\renewcommand{\EnvelopeThreeVarRmseIQ}{6.15}
\renewcommand{\EnvelopeThreeVarRmseMC}{8.05}
\renewcommand{\EnvelopeThreeVarSkillIQ}{0.46}
\renewcommand{\EnvelopeThreeVarSkillMC}{0.30}


\title{Supporting Material: Artifact-Level Detail for Rank Instability and
Representation Repair in Automatic Modulation Classification under Structured
Interference}
\author{Anonymous Author(s)}

\begin{document}
\maketitle

\begin{abstract}
This is optional supporting material. No result required to support the
manuscript's main claims appears only in this document. It provides complete
cell construction, aggregation definitions, per-seed/SNR/class results,
architecture details, frozen-control analyses, convergence histories,
prospective-freeze provenance, and artifact paths for results summarized in the
manuscript. The three prospectively frozen severe-held stages---the transfer
check, the sidecar repair, and the frozen-inference attribution controls---are
stated in the manuscript and expanded here only at artifact level; none
changes the original sealed evidence class. A reviewer who wants to audit the
construction of the evaluation cells, the full capacity-control sweep, the
per-SNR comparator-fidelity curves, the covariate ladder and selector audits,
the leave-$-15$-dB-out extrapolation and paired overlap-increment tests, the
comparator convergence histories, the sidecar architecture, the severe-held
construction and control breakdowns, or the seed-level robustness can do so
directly.
\end{abstract}

\section{Scope and How to Read This Document}
The study asks why compact spectral and high-capacity I/Q AMC representations
exhibit condition-dependent rank instability, how far a campaign-internal
empirical envelope transports, and whether a joint representation-and-fusion
intervention can repair the degradation exposed by independent severe-domain
validation. All
evidence is from source-disjoint 3GPP TR~38.901 TDL-profile simulation.

Because this document is standalone, the identifiers used throughout are
restated here.
\begin{itemize}
  \item \textbf{A0} is the direct spectral reference backbone (8{,}458
  parameters) and \textbf{A5} is the sealed complete configuration
  (39{,}500 parameters), which adds a three-route spectral allocation, a
  simulation-component teacher, multi-task and contrastive objectives, a
  residual bypass, and a small route-orthogonality regularizer.
  \item \textbf{S}, \textbf{M}, and \textbf{L} are prospective spectral
  capacity tiers. S \emph{is} the sealed A5 configuration; M and L widen only
  the spectral-channel, embedding, and environment dimensions, to
  99{,}596 and 233{,}244 parameters.
  \item \textbf{sidecar} is a prospective received-I/Q repair branch that
  raises the total model size from 39{,}500 to 46{,}794 parameters.
  \item \textbf{IQFormer-inspired} and \textbf{MCLDNN} are
  architecture-faithful unified-retraining references, not exact reproductions
  of every original training recipe.
  \item The \textbf{target-support jammer overlap} $o$ is the fraction of the
  jammer's total STFT power that falls inside the target's 99\%-energy
  time--frequency support. It is not an occupied-bandwidth fraction.
\end{itemize}

Four evidence levels are kept apart. The sealed family contains the original
confirmatory contrasts. Campaign-internal SNR--SIR, envelope, fusion, and
capacity analyses are exploratory. Probe, coverage, route-gating, and sidecar
experiments are prospective diagnosis or intervention evidence. The severe-held
transfer check, the sidecar repair, and the frozen-inference attribution
controls are three sequential prospectively frozen post-campaign diagnostic
stages on one independently generated held set; their freeze provenance is in
Section~\ref{sec:preregp}. None changes the original sealed family.

\section{Construction of the Evaluation Cells}
The envelope is fitted on aggregate cells rather than on individual windows. A
cell is one (split, SIR stratum, overlap quartile) triple that retains at least
150 windows and all ten modulation classes. Overlap quartile edges are computed
within each split. Table~\ref{tab:cells} gives the resulting construction; the
fit total of 32 and the primary held total of 80 are both reconstructible from
its rows.

\begin{table}[t]
\centering
\caption{Cell construction for the campaign-diagnostic envelope. ``Windows per cell'' is
the range of source-disjoint test windows contributing to a cell.}
\label{tab:cells}
\footnotesize
\setlength{\tabcolsep}{3pt}
\begin{tabular}{llrrr}
\toprule
Split & Role & SIR strata (dB) & Cells & Windows per cell\\
\midrule
\texttt{id\_test} & fit & $-$10, $-$5, 0, $+$5, $+$10 & 16 & 230--1053\\
\texttt{hard\_interference} & fit & $-$15, $-$10, $-$5, 0 & 16 & 290--343\\
\texttt{unseen\_jammer} & held & $-$10, $-$5, 0, $+$5, $+$10 & 20 & 229--276\\
\texttt{unseen\_speed} & held & $-$10, $-$5, 0, $+$5, $+$10 & 20 & 228--277\\
\texttt{heldout\_channel} & held & $-$10, $-$5, 0, $+$5, $+$10 & 20 & 216--279\\
\texttt{combined\_ood} & held & $-$10, $-$5, 0, $+$5, $+$10 & 20 & 223--270\\
\texttt{clean\_retention} & excluded & $0$ & 1 & 5000\\
\midrule
\multicolumn{3}{l}{Fit total} & 32 & \\
\multicolumn{3}{l}{Primary held total} & 80 & \\
\bottomrule
\end{tabular}
\end{table}

Two features of Table~\ref{tab:cells} are worth stating plainly. First, the
\texttt{id\_test} split contributes 16 cells rather than 20 because its lowest
overlap quartile is populated almost entirely by jammer-free windows, which
carry $o=0$ and the SIR $=0$~dB cache sentinel; that bin clears the 150-window
floor only in the $0$~dB stratum. Second, the single
\texttt{clean\_retention} cell is excluded from the primary envelope because a
jammer-free window has no physical finite SIR, which is also why its window
count is the full split rather than a quartile of it.

\section{Campaign-Diagnostic Envelope Skill by Held Regime}
Table~\ref{tab:regime} reports the held-regime skill for both strong
references, including two diagnostics that are easy to omit and easy to
misread. ``Const.'' is the fit-domain constant predictor: the mean gain over
the 32 fitting cells, applied unchanged to the held cells, with
$R^2_{\mathrm{skill}}=1-\mathrm{MSE}_{\mathrm{envelope}}/\mathrm{MSE}_{\mathrm{const}}$.
``Oracle'' is instead the mean of the held cells themselves (pooled row) or of
the regime's own held cells (regime rows), which no predictor could know in
advance; it is reported as a hindsight reference, not as evidence that the
envelope fails to resolve within-regime structure.

\begin{table}[t]
\centering
\caption{Held-regime envelope skill. RMSE values in \pp.}
\label{tab:regime}
\scriptsize
\setlength{\tabcolsep}{2.4pt}
\begin{tabular}{llrrrrrr}
\toprule
Reference & Held regime & Cells & Env. & Const. & $R^2_{\mathrm{skill}}$ & Oracle & Sign\\
\midrule
IQFormer-insp. & unseen speed & 20 & 3.94 & 7.87 & 0.75 & 6.71 & 1.000\\
IQFormer-insp. & held-out channel & 20 & 4.86 & 7.59 & 0.59 & 6.95 & 1.000\\
IQFormer-insp. & combined OOD & 20 & 7.50 & 9.23 & 0.34 & 4.24 & 0.950\\
IQFormer-insp. & unseen jammer & 20 & 7.74 & 8.67 & 0.20 & 3.83 & 0.950\\
IQFormer-insp. & pooled (80) & 80 & 6.23 & 8.37 & 0.45 & 6.04 & 0.975\\
\midrule
MCLDNN & unseen speed & 20 & 4.66 & 8.88 & 0.72 & 7.78 & 1.000\\
MCLDNN & held-out channel & 20 & 5.41 & 9.06 & 0.64 & 8.41 & 1.000\\
MCLDNN & combined OOD & 20 & 10.23 & 10.28 & 0.01 & 2.71 & 0.900\\
MCLDNN & unseen jammer & 20 & 10.34 & 10.24 & -0.02 & 3.23 & 0.900\\
MCLDNN & pooled (80) & 80 & 8.10 & 9.63 & 0.29 & 6.81 & 0.950\\
\bottomrule
\end{tabular}
\end{table}

The heterogeneity is real and is not smoothed away. Against
IQFormer-inspired, the envelope predicts the gain level well in the
unseen-speed and held-out-channel regimes ($R^2_{\mathrm{skill}}=0.75$ and
$0.59$) and much less well into the unseen-jammer and combined-OOD regimes
($0.20$ and $0.34$); against MCLDNN the unseen-jammer and combined-OOD regimes
are at or below the constant baseline. The regime-wise oracle comparison is
likewise heterogeneous: on the unseen-speed and held-out-channel shifts the
envelope substantially beats the corresponding regime-specific constant
(3.94 vs 6.71 and 4.86 vs 6.95 pp against IQFormer-inspired; 4.66 vs 7.78
and 5.41 vs 8.41 pp against MCLDNN), while on the unseen-jammer and
combined-OOD shifts the regime-specific constant is better. The pooled oracle
(6.04 and 6.81 pp) is lower than the pooled envelope RMSE because it
uses the held-domain mean in hindsight; a pooled constant still bears the
between-regime variation and the pooled comparison should not be read as the
envelope lacking within-regime structure. The demonstrated property is
therefore regime-dependent transport and within-regime structure, not uniform
held-domain superiority.

\begin{figure}[t]
\centering
\includegraphics[width=0.92\columnwidth]{../paper_figures/outputs_v46/fig3_campaign_diagnostic_envelope.pdf}
\caption{Exploratory campaign-diagnostic overlap--SIR envelope against
IQFormer-inspired (moved from the manuscript's figure set). Panel~(a) uses
the target-support jammer overlap, not occupied bandwidth: filled markers are
the 32 fit cells and open markers the 80 held cells. Panel~(b) compares
predicted and observed held-cell gain. RMSE is
\EnvelopeIqformerHoldoutRmse\ \pp\ versus \EnvelopeConstantRmseIQ\ \pp\ for
the frozen fit-domain constant
($R^2_{\mathrm{skill}}=\EnvelopeSkillRTwoIQ$). The surface has limited
held-regime skill inside the frozen campaign and does not transport to the
independent severe-held regimes (Section~\ref{sec:transfer}).}
\label{fig:envelope_supp}
\end{figure}

\section{Clean-Sentinel Fitting-Cell Sensitivity}
The primary held envelope excludes jammer-free clean retention because a
jammer-free window has no physical finite SIR. The fitting domain, however,
retains one aggregate dominated by jammer-free sentinel windows, the
\texttt{id\_test}/SIR$=0$/lowest-overlap cell (mean overlap
$\approx10^{-11}$, 1053 windows, SIR taken from the cache sentinel). To
close this definition asymmetry, we exclude exactly that aggregate, refit
the two-variable envelope on the remaining 31 cells, and re-score the
unchanged 80 finite-SIR interfered held cells; no model prediction changes
and nothing is retrained.

\begin{table}[t]
\centering
\caption{Clean-sentinel fitting-cell sensitivity. The 31-cell fit excludes
the single \texttt{id\_test}/SIR$=0$/lowest-overlap aggregate dominated by
jammer-free sentinel windows; the held evaluation is the unchanged 80
finite-SIR interfered cells. RMSE values in \pp; BE denotes the break-even
overlap.}
\label{tab:sentinel}
\scriptsize
\setlength{\tabcolsep}{2.0pt}
\begin{tabular}{lrrrrrrrr}
\toprule
Reference & Fit & Held & Const. & $R^2_{\mathrm{skill}}$ & Overlap & SIR & BE $o$ & BE $o$\\
& cells & RMSE & RMSE & & coeff. & coeff. & @ $-15$ & @ $-10$\\
\midrule
IQFormer-insp. & 32 & 6.23 & 8.37 & 0.45 & 14.38 & $-1.036$ & 0.11 & 0.47\\
IQFormer-insp. & 31 & 6.42 & 8.28 & 0.40 & 15.90 & $-1.066$ & 0.14 & 0.47\\
\midrule
MCLDNN & 32 & 8.10 & 9.63 & 0.29 & 19.16 & $-0.894$ & 0.06 & 0.29\\
MCLDNN & 31 & 8.20 & 9.59 & 0.27 & 20.55 & $-0.922$ & 0.08 & 0.31\\
\bottomrule
\end{tabular}
\end{table}

The excluded aggregate is the \texttt{id\_test}/SIR$=0$/lowest-overlap cell
dominated by jammer-free sentinel windows. Excluding it moves the held RMSE
by at most $0.20$~\pp\ and $R^2_{\mathrm{skill}}$ by at most $0.05$, leaves
both coefficient signs and the break-even ordering (required overlap rising
as interference becomes less severe) unchanged, and therefore leaves the
envelope conclusions materially unchanged; the 32-cell primary fit is
retained.

\section{Capacity Control across All Splits}
The capacity ladder isolates width from representation: S, M, and L differ only
in three width numbers and share the routes, teacher, objective, and training
recipe. All values in Tables~\ref{tab:capabs} and~\ref{tab:capdiff} are
computed on the same matched five-seed subset, and must not be mixed
numerically with ten-seed results reported elsewhere.

\begin{table}[t]
\centering
\caption{Absolute macro-F1 on the matched five-seed subset. These split-level
five-seed macro-F1 values should not be numerically differenced against the
ten-seed cell-mean gains in the manuscript envelope table.}
\label{tab:capabs}
\footnotesize
\setlength{\tabcolsep}{4pt}
\begin{tabular}{lrrrrr}
\toprule
Split & S & M & L & IQFormer-insp. & MCLDNN\\
\midrule
hard interference & 0.436 & 0.455 & 0.468 & 0.496 & 0.454\\
hard, SIR $=-15$ dB & 0.353 & 0.384 & 0.403 & 0.262 & 0.255\\
in distribution & 0.563 & 0.584 & 0.588 & 0.696 & 0.636\\
unseen jammer & 0.456 & 0.462 & 0.466 & 0.612 & 0.584\\
unseen speed & 0.487 & 0.507 & 0.510 & 0.626 & 0.582\\
held-out channel & 0.488 & 0.510 & 0.515 & 0.613 & 0.567\\
combined OOD & 0.450 & 0.462 & 0.465 & 0.621 & 0.587\\
clean retention & 0.500 & 0.511 & 0.512 & 0.693 & 0.643\\
\bottomrule
\end{tabular}
\end{table}

\begin{table}[t]
\centering
\caption{Paired capacity contrasts on the matched five-seed subset, in \pp\
with 95\% intervals.}
\label{tab:capdiff}
\footnotesize
\setlength{\tabcolsep}{4pt}
\begin{tabular}{lrrr}
\toprule
Contrast (hard interference unless noted) & Diff. & Low & High\\
\midrule
M minus S & $1.94$ & $1.07$ & $2.80$\\
L minus S & $3.18$ & $2.44$ & $3.96$\\
L minus IQFormer-insp. & $-2.84$ & $-4.30$ & $-1.31$\\
L minus MCLDNN & $1.37$ & $0.03$ & $2.72$\\
L minus S, SIR $=-15$ dB & $4.96$ & $3.48$ & $6.48$\\
L minus IQFormer-insp., SIR $=-15$ dB & $14.07$ & $9.55$ & $18.51$\\
L minus IQFormer-insp., clean retention & $-18.14$ & $-19.34$ & $-17.00$\\
\bottomrule
\end{tabular}
\end{table}

The pattern is consistent across splits. Width raises the average
monotonically, and it helps most where the compact model was already
competitive: at SIR $=-15$~dB the widest tier gains $4.96$~\pp\ over S and
stays $14.07$~\pp\ ahead of IQFormer-inspired. It does not close the gap that
matters for the boundary. On the marginal hard split L is still $2.84$~\pp\
behind IQFormer-inspired, and on clean retention it is $18.14$~\pp\ behind,
a deficit essentially untouched by tripling the width from M to L
($0.10$~\pp\ of movement). This is the arithmetic behind the manuscript's
statement that capacity is not a substitute for the representation change.

\section{Prospective Tier-2 Interventions}
The Tier-2 chain was designed after inspecting the frozen evidence and sits
outside the sealed family. Table~\ref{tab:tier2s} restates its three arms with
the intervals that the compact form in the manuscript omits.

\begin{table}[t]
\centering
\caption{Prospective Tier-2 diagnosis, repair, and negative control, relative
to A5, in \pp\ with 95\% intervals. Exploratory throughout.}
\label{tab:tier2s}
\footnotesize
\setlength{\tabcolsep}{2.6pt}
\begin{tabular}{lrrr}
\toprule
Intervention & Params & Clean diff. [95\% CI] & Hard diff. [95\% CI]\\
\midrule
Per-class clean coverage & 39.5k & $0.15$ $[-0.31,0.63]$ & $-0.30$ $[-0.91,0.27]$\\
Presence-gated route & 39.5k & $-0.30$ $[-0.66,0.08]$ & $-0.18$ $[-0.72,0.35]$\\
Received-I/Q sidecar & 46.8k & $11.77$ $[11.06,12.49]$ & $4.37$ $[3.44,5.31]$\\
\bottomrule
\end{tabular}
\end{table}

The ordering carries the argument. Adding jammer-free windows for every class
moves clean macro-F1 by $0.15$~\pp, so exposure alone is not the explanation.
Gating the modulation route on predicted interference presence moves nothing
either, and the learned gate saturates near the modulated route in both clean
and hard conditions ($99.994\%$ and $99.991\%$), so it functions as a negative
control rather than as a demonstration that a mask fallback was learned. Only
the full received-I/Q sidecar configuration moves both clean and hard
conditions with intervals clear of zero. The later frozen-inference controls
in Section~\ref{sec:controls} show that the effect depends on received-I/Q
information but also requires the co-trained fusion/classification path, so it
should not be interpreted as an isolated I/Q-information effect.

\section{Seed-Level Robustness}
The sealed intervals treat algorithm seeds as fixed blocks. Treating them
instead as a random effect and resampling seeds directly gives
$[4.02,5.19]$~\pp\ for the A5--A0 whole-configuration contrast,
$[-0.15,0.42]$~\pp\ for the teacher-form contrast, $[-0.15,1.04]$~\pp\ for the
route-count contrast, and $[-0.19,0.62]$~\pp\ for the residual-free A7 variant
against A5. The qualitative reading is unchanged under either treatment: the
whole-configuration effect is resolved, the isolated component effects are not.

\section{Comparator Fidelity on a Public Benchmark}
This check exists to establish that the two I/Q comparators are implemented
competently, not to validate the simulator-domain envelope; a public benchmark
cannot do the latter, because it does not expose the separately tracked target
and jammer components the overlap covariate requires. The protocol follows the
benchmark result table released with the IQFormer reference implementation:
the 11-class RML2016.10a set, the full $-20$ to $18$~dB grid, per-stratum
$600/200/200$ splits seeded at $233$, three algorithm seeds, no augmentation
and no input-length change. The public-benchmark fidelity runs use a separate
recipe following the benchmark protocol---a $60$-epoch budget for
IQFormer-inspired and a $200$-epoch budget for MCLDNN---and should not be used
to infer convergence sufficiency of the $30$-epoch main campaign.

\begin{table}[t]
\centering
\caption{Per-SNR accuracy on RML2016.10a over three seeds. Brackets are the
seed minimum and maximum, not confidence intervals.}
\label{tab:e1s}
\footnotesize
\setlength{\tabcolsep}{3.4pt}
\begin{tabular}{rrrrr}
\toprule
SNR (dB) & MCLDNN mean & [min, max] & IQFormer-insp. mean & [min, max]\\
\midrule
$-20$ & 12.1\% & $[12.0,12.2]$ & 14.9\% & $[14.5,15.2]$\\
$-18$ & 13.1\% & $[12.3,13.9]$ & 18.1\% & $[16.8,19.0]$\\
$-16$ & 14.8\% & $[14.0,15.3]$ & 22.3\% & $[21.7,23.0]$\\
$-14$ & 17.9\% & $[16.8,18.9]$ & 27.2\% & $[26.6,28.3]$\\
$-12$ & 22.0\% & $[20.4,23.3]$ & 33.3\% & $[31.5,34.6]$\\
$-10$ & 26.0\% & $[25.0,26.7]$ & 37.0\% & $[36.5,37.8]$\\
$-8$ & 32.5\% & $[31.1,33.3]$ & 41.4\% & $[41.1,41.5]$\\
$-6$ & 44.0\% & $[42.8,45.7]$ & 49.8\% & $[49.5,50.3]$\\
$-4$ & 58.8\% & $[58.3,59.9]$ & 63.3\% & $[63.0,63.5]$\\
$-2$ & 74.7\% & $[74.3,75.0]$ & 77.3\% & $[76.1,77.9]$\\
$0$ & 83.4\% & $[82.9,84.1]$ & 86.1\% & $[85.7,86.7]$\\
$2$ & 85.1\% & $[84.7,85.6]$ & 89.4\% & $[88.6,90.5]$\\
$4$ & 85.5\% & $[84.7,87.1]$ & 91.6\% & $[90.6,93.0]$\\
$6$ & 86.1\% & $[85.7,86.8]$ & 93.4\% & $[93.0,94.1]$\\
$8$ & 86.4\% & $[84.9,88.6]$ & 93.3\% & $[92.8,94.2]$\\
$10$ & 85.3\% & $[84.0,87.1]$ & 93.0\% & $[92.0,94.4]$\\
$12$ & 86.3\% & $[85.1,88.4]$ & 93.8\% & $[93.1,94.7]$\\
$14$ & 86.0\% & $[85.0,87.9]$ & 93.4\% & $[92.9,93.9]$\\
$16$ & 85.5\% & $[84.5,87.1]$ & 92.9\% & $[92.1,93.7]$\\
$18$ & 85.2\% & $[84.3,86.8]$ & 92.9\% & $[92.3,93.9]$\\
\midrule
All SNR & 58.55\% & $[58.1,59.4]$ & 65.23\% & $[64.8,65.8]$\\
Reference release & 62.05\% & --- & 64.19\% & ---\\
\bottomrule
\end{tabular}
\end{table}

Table~\ref{tab:e1s} makes the shape of the agreement explicit. Both
implementations reproduce the reference release at the level of the all-SNR
range, and IQFormer-inspired matches it closely ($+1.04$~\pp), while MCLDNN
remains $3.50$~\pp\ lower but within the same broad range. Agreement is not
pointwise: individual SNR bins deviate by up to about $8$~\pp, which is why
this is reported as range-level implementation fidelity rather than
reproduction. No fidelity threshold was declared before these runs, so no
pass/fail band is claimed here.

\section{Comparator Convergence Audit}
All models were trained under one shared 30-epoch budget with the same
cosine-scheduled learning rate and an 8-epoch early-stopping rule; the
reported checkpoint is selected solely by validation modulation
cross-entropy. Table~\ref{tab:convergence} audits the sealed training
histories (selected epoch and completed epochs over the ten algorithm seeds)
to check whether any comparator's comparison is distorted by an immediately
binding training cap.

\begin{table}[t]
\centering
\caption{Convergence audit. ``At epoch 30'' counts seeds whose training ran the
full shared budget; ``Early-stopped'' counts seeds stopped earlier by the
patience rule. Panel (a) is sealed; panel (b) reads the post-campaign histories.}
\label{tab:convergence}
\footnotesize
\setlength{\tabcolsep}{5pt}
\textbf{(a) Sealed comparator histories}\\[2pt]
\begin{tabular}{lrrrr}
\toprule
Model & Median best epoch & Range & Seeds at epoch 30 & Early-stopped seeds\\
\midrule
MCLDNN & 13 & 11--16 & 0/10 & 10/10\\
IQFormer-insp. & 27 & 24--30 & 10/10 & 0/10\\
CSSL & 29 & 26--30 & 10/10 & 0/10\\
A5 & 29.5 & 27--30 & 10/10 & 0/10\\
A0 & 30 & 28--30 & 10/10 & 0/10\\
\bottomrule
\end{tabular}

\vspace{5pt}
\textbf{(b) Post-campaign capacity and repair histories}\\[2pt]
\begin{tabular}{lrrrr}
\toprule
Model & Median best epoch & Range & Seeds at epoch 30 & Early-stopped seeds\\
\midrule
Capacity M & 29 & 27--30 & 5/5 & 0/5\\
Capacity L & 27 & 24--30 & 5/5 & 0/5\\
I/Q sidecar & 30 & 29--30 & 10/10 & 0/10\\
\bottomrule
\end{tabular}
\end{table}

The budget binds asymmetrically. IQFormer-inspired and CSSL run the full
30-epoch budget in all ten seeds (median selected epoch 27 and 29), so their
reported performance should be interpreted as potentially budget-limited under
this shared 30-epoch protocol and the shared epoch budget may favor
faster-converging
models; the manuscript states this caveat in its Limitations. MCLDNN
early-stops in every seed with best epochs well inside the budget (median 13,
range 11--16), so the severe-region comparisons against MCLDNN are not
explained by a binding cap. A5 itself also trains to the budget in all ten
seeds, and the direct spectral reference A0 likewise runs to the budget in
all ten seeds (median selected epoch 30, range 28--30), so the cap concern
is not one-sided against the comparators. No model
was retrained for this audit. The post-campaign M, L, and sidecar runs also
complete the full budget in every seed, with median selected epochs 29, 27,
and 30, respectively; this records their optimization histories without
promoting the diagnostic stages into the sealed family.

\section{Sidecar Architecture Detail}
Table~\ref{tab:sidecar} lists the received-I/Q sidecar layer by layer so the
reported parameter increment is fully reconstructible. All convolution layers
are bias-free; SiLU activations (parameter-free) follow each normalization.
The branch consumes the received two-channel mixture at the full sample rate,
uses one depthwise-separable stack with aggressive time pooling
(mean/std/max over the window), and fuses the resulting embedding residually
with the spectral embedding before a side classification head. The fused
embedding passes through the 10-class linear head of layer~13, and these
logits are added residually to the base spectral classifier logits; the side
head receives no separate auxiliary supervision, and inference consumes only
the received mixture. The parameter
counts were verified directly against the sealed checkpoint state dict.

\begin{table}[t]
\centering
\caption{Sidecar layer structure and parameter count. All Conv1D layers use
stride 1 and no bias; DW denotes depthwise (groups = in channels) and PW
pointwise ($1\times1$) convolution. Total model size is $39{,}500+7{,}294=
46{,}794$ parameters.}
\label{tab:sidecar}
\footnotesize
\setlength{\tabcolsep}{4pt}
\begin{tabular}{llccrr}
\toprule
Layer & Type & Kernel & Stride & In$\to$Out ch. & Params\\
\midrule
1 & Conv1D & 9 & 1 & 2$\to$12 & 216\\
2 & GroupNorm (3 groups) & --- & --- & 12 & 24\\
3 & DW Conv1D & 7 & 1 & 12$\to$12 & 84\\
4 & PW Conv1D & 1 & 1 & 12$\to$24 & 288\\
5 & GroupNorm (6 groups) & --- & --- & 24 & 48\\
6 & DW Conv1D & 5 & 1 & 24$\to$24 & 120\\
7 & PW Conv1D & 1 & 1 & 24$\to$24 & 576\\
8 & GroupNorm (6 groups) & --- & --- & 24 & 48\\
9 & mean/std/max pool + Linear & --- & --- & 72$\to$24 & 1{,}752\\
10 & LayerNorm & --- & --- & 24 & 48\\
11 & Fusion Linear & --- & --- & 72$\to$48 & 3{,}504\\
12 & Fusion LayerNorm & --- & --- & 48 & 96\\
13 & Classifier Linear & --- & --- & 48$\to$10 & 490\\
\midrule
\multicolumn{5}{l}{Total sidecar increment} & 7{,}294\\
\bottomrule
\end{tabular}
\end{table}

\section{Diagnostic Envelope Covariate and Selector Audits}
Two zero-compute audits recompute, on the same 32 fitting cells and the same
80 finite-SIR interfered held cells as the primary envelope, (i) a covariate
ladder that isolates the skill of each physical covariate and (ii) a
diagnostic selector evaluation. No model is retrained and no prediction
changes; both are exploratory.

Table~\ref{tab:ladder} gives the ladder. All models are fitted on the 32
fitting cells by least squares and scored unchanged on the 80 held cells;
$R^2_{\mathrm{skill}}$ uses the same fit-domain constant baseline as the
manuscript. SIR is the stronger single covariate against IQFormer-inspired,
whereas SIR and overlap have comparable single-covariate skill against
MCLDNN; overlap adds a consistent increment beyond SIR against both
references, while SIR adds more beyond overlap against IQFormer-inspired
than the reverse. The statistical status of that increment is assessed by the
paired bootstrap in Section~\ref{sec:s4tests}.

\begin{table}[t]
\centering
\caption{Covariate ladder on the shared 32 fitting / 80 held cells. RMSE in
\pp. Increments: overlap$|$SIR reduces held RMSE by $0.61$~\pp\ against both
references; SIR$|$overlap by $1.23$~\pp\ against IQFormer-inspired and
$0.58$~\pp\ against MCLDNN.}
\label{tab:ladder}
\footnotesize
\setlength{\tabcolsep}{4pt}
\begin{tabular}{lrrrr}
\toprule
Model & IQ RMSE & IQ $R^2_{\mathrm{skill}}$ & MC RMSE & MC $R^2_{\mathrm{skill}}$\\
\midrule
Constant & 8.37 & 0.00 & 9.63 & 0.00\\
SIR only & 6.84 & 0.33 & 8.71 & 0.18\\
Overlap only & 7.46 & 0.20 & 8.68 & 0.19\\
Overlap + SIR & 6.23 & 0.45 & 8.10 & 0.29\\
Overlap + SIR + SNR & 6.15 & 0.46 & 8.05 & 0.30\\
\bottomrule
\end{tabular}
\end{table}

Table~\ref{tab:selector} evaluates the envelope sign as a diagnostic rule on
the 80 held cells. Each sign selector uses envelope coefficients fitted only
on the 32 fitting cells and applied to held cells by covariate value; the
oracle is a hindsight per-cell best and is not achievable. Realized score is
the mean over cells of the selected model's actual macro-F1; MACs use the
sealed complexity table values (41.8~M for A5, 355.6~M for IQFormer-inspired).

\begin{table}[t]
\centering
\caption{Selector evaluation on the 80 held cells. Regret is against the
hindsight oracle; expected MACs average the selected model's cost over cells.
These selector statistics are computed from the ten-seed held-cell aggregates
and should not be numerically differenced against the matched-five-seed
split-level macro-F1 values in Table~\ref{tab:capabs}.}
\label{tab:selector}
\footnotesize
\setlength{\tabcolsep}{3pt}
\begin{tabular}{lrrrrr}
\toprule
Selector & Realized F1 & Regret (pp) & A5 cells & MACs (M) & MAC saving\\
\midrule
Always A5 & 0.4488 & 14.19 & 80/80 & 41.8 & 88.2\%\\
Always IQFormer & 0.5894 & 0.14 & 0/80 & 355.6 & 0\%\\
SIR-only sign rule & 0.5894 & 0.14 & 0/80 & 355.6 & 0\%\\
Overlap+SIR sign rule & 0.5872 & 0.36 & 5/80 & 336.0 & 5.5\%\\
Oracle (hindsight) & 0.5907 & 0 & 3/80 & 343.8 & 3.3\%\\
\bottomrule
\end{tabular}
\end{table}

The reading is deliberately plain. Inside the held finite-SIR regimes the
large I/Q reference leads in almost every cell, so no selector recovers much
against always using it, and the envelope's demonstrated value is boundary
localization---identifying the few severe, high-overlap cells where the
compact model is competitive---rather than a large expected compute saving in
this campaign. The overlap+SIR rule is the only rule that ever selects the
compact model (5 of 80 cells), at a regret of $0.36$~\pp\ against the
hindsight oracle; the SIR-only rule never selects it. Computed directly from
the held-cell artifacts, its confusion matrix against the hindsight-optimal
indicator is TP$=3$, FP$=2$, FN$=0$, TN$=75$: it identifies all three
compact-favoring cells (recall $3/3$) with two additional selections
(precision $3/5$), so its held-domain value is boundary localization without
missed positive cells, while always using IQFormer-inspired still attains the
higher average realized macro-F1. This is the honest
selector consequence of a regime structure in which the compact model's
advantage is concentrated in a small corner of the plane.

\section{Leave-\texorpdfstring{$-15$}{-15}-dB-Out Extrapolation and Paired Overlap-Increment Tests}
\label{sec:s4tests}
Two further zero-compute audits close the remaining envelope questions. Both
reuse the frozen cell table; no model is retrained and no prediction changes.

The 80 held cells contain no SIR $=-15$~dB cell: that stratum occurs only in
the \texttt{hard\_interference} fitting split. To probe extrapolation toward
the most severe level, we refit the two-variable envelope on the 28 fitting
cells excluding the four $-15$-dB cells and predict those four cells
out-of-sample (a leave-one-SIR-stratum-out test, not an external validation).
Table~\ref{tab:leave15} compares the envelope with a 28-cell constant and a
28-cell SIR-only fit.

\begin{table}[t]
\centering
\caption{Leave-$-15$-dB-out test. All predictors are fitted on the 28
fitting cells excluding the four \texttt{hard\_interference} SIR $=-15$~dB
cells and scored on those four cells only. RMSE and MAE in \pp; sign
accuracy is the fraction of the four cells whose predicted gain sign matches
the observed sign.}
\label{tab:leave15}
\footnotesize
\setlength{\tabcolsep}{4pt}
\begin{tabular}{llrrr}
\toprule
Reference & Predictor & RMSE & MAE & Sign acc.\\
\midrule
IQFormer-insp. & Constant & 21.34 & 21.26 & 0/4\\
IQFormer-insp. & SIR only & 12.15 & 12.02 & 0/4\\
IQFormer-insp. & Overlap + SIR & 12.69 & 12.01 & 1/4\\
\midrule
MCLDNN & Constant & 16.43 & 15.83 & 0/4\\
MCLDNN & SIR only & 7.87 & 6.54 & 4/4\\
MCLDNN & Overlap + SIR & 6.84 & 6.52 & 2/4\\
\bottomrule
\end{tabular}
\end{table}

The outcome is reference-dependent, and Table~\ref{tab:leave15} must be read
on both skill dimensions. For MCLDNN, the full overlap+SIR surface reduces
RMSE relative to the constant baseline ($6.84$ versus $16.43$~\pp) but
predicts the correct sign in only two of four held-out cells, whereas the
SIR-only surface preserves all four signs ($7.87$~\pp\ RMSE). For
IQFormer-inspired, the full model is no better than SIR-only ($12.69$ versus
$12.15$~\pp\ RMSE; sign accuracy $1/4$ versus $0/4$). The
leave-one-SIR-stratum-out test therefore does not confirm a transportable
overlap increment: the full surface does not preserve the severe-region sign
structure that the SIR-only surface retains. The $-15$-dB break-even overlaps
remain conditional extrapolations of the fitted envelope rather than
independently held-regime-validated thresholds.

For the incremental value of overlap beyond SIR on the 80 held cells, we use
a block-paired bootstrap. The paired statistic is the per-cell
squared-error difference
$d_i=e^2_{\mathrm{SIR\text{-}only},i}-e^2_{\mathrm{overlap+SIR},i}$, so
$\bar d>0$ means adding overlap lowers MSE. The 80 held cells form 20
held-regime$\times$SIR blocks of four overlap quartiles each; each draw
resamples the 20 blocks with replacement ($B=10{,}000$, seed $20260820$)
and keeps all quartiles inside a sampled block.

\begin{table}[t]
\centering
\caption{Block-paired bootstrap of the overlap increment beyond SIR on the
80 held cells. $\Delta$MSE is the paired squared-error reduction
(SIR-only minus overlap+SIR), in $\mathrm{pp}^2$; intervals are 2.5/97.5
percentiles over $10{,}000$ block resamples.}
\label{tab:bootpair}
\footnotesize
\setlength{\tabcolsep}{4pt}
\begin{tabular}{lrrrr}
\toprule
Reference & Mean $\Delta$MSE & 95\% CI & $P(\Delta\mathrm{MSE}>0)$ & RMSE red. (pp)\\
\midrule
IQFormer-insp. & 7.97 & $[-2.06, 18.18]$ & 0.94 & 0.61\\
MCLDNN & 10.18 & $[-8.29, 29.84]$ & 0.85 & 0.61\\
\bottomrule
\end{tabular}
\end{table}

Both references show positive point estimates and a majority of positive
draws, but both percentile intervals include zero; overlap therefore provides
a consistent point-estimate refinement beyond SIR rather than a separately
resolved predictor, and the manuscript reports it as an exploratory
refinement.

Finally, applying the 28-cell sign rule to the four excluded $-15$-dB cells
with the same ten-seed cell-level aggregates and MAC cost model selects A5
in only one of the four hindsight-A5-favoring cells (sign accuracy $1/4$;
realized cell-mean macro-F1 $0.285$ versus the $0.361$ hindsight oracle on
those four cells). The cross-model selector therefore transfers poorly to
the excluded severe stratum, which is the selector-side counterpart of the
extrapolation caveat above.

\section{Independent Severe-Held Validation: Transfer Boundary of the Corner
Advantage and the Diagnostic Envelope}
\label{sec:transfer}
After the campaign was frozen, a prospectively frozen post-campaign robustness
validation asked whether the severe-corner advantage and the overlap--SIR
envelope survive on data generated outside the frozen campaign. The
independent severe-held set fixes SIR at $-15$~dB and uses two
source-disjoint regimes with different operating-regime factors: regime R1
keeps the training channel profiles (TDL-A/C/D) and combines them with held
jammer families (pulse and OFDM-like emitters) at a held terminal speed of
$180$~km/h, while regime R2 combines held channel profiles (TDL-B/E) with the
same held jammer families at $250$~km/h. Each regime contains
\SevereHeldSourcesPerRegime\ fresh sources disjoint from the campaign source
pool (\SevereHeldSourcesTotal\ new sources against the campaign's 152,000; the disjointness is
asserted programmatically at cache build), evaluated on the eight-point SNR
grid from $-10$ to $18$~dB at view $0$. The set is stored under
\path{standards/cache_s5_severe_held_v1/} with a manifest recording the
cache digest. All frozen checkpoints were evaluated unchanged; no model was
retrained, tuned, or selected against this set. The validation is
exploratory and post-campaign relative to the sealed family throughout.

Two aggregation conventions appear in this document, and every reported
number is labelled by its convention. \emph{Window-pooled} macro-F1 pools all
predictions first and then computes macro-F1; it is used for the transfer
gains in Table~\ref{tab:s5a}. The \emph{seed-wise} convention computes
macro-F1 per seed and per regime first and then averages, with pooled the
equally weighted regime mean; it is used for the repair contrast and levels
and for all frozen-inference controls, and it is the manuscript's headline
convention because it supports paired seed inference. Because macro-F1 is
nonlinear, the two conventions need not agree numerically
($\SevereTransferGainIqformer$ versus $\SevereSeedMeanGapIqformer$ \pp\
against IQFormer-inspired, $\SevereTransferGainMcldnn$ versus
$\SevereSeedMeanGapMcldnn$ \pp\ against MCLDNN for the A5 gap on this held
set), and they are never mixed within a single contrast.

Table~\ref{tab:s5a} reports the frozen transfer verdicts. The pooled
A5 gain is negative against both comparators, the envelope's sign accuracy
is near chance (and below it against MCLDNN), a SIR-only or constant
predictor beats the full overlap--SIR envelope on RMSE against both
references, and the overlap--SIR cell selector recalls only
\SevereTransferSelectorRecall\ of the truly A5-favoring cells. The
frozen verdict is therefore a transfer failure: neither the
severe-corner advantage nor the $-15$-dB envelope transport reproduces on
independently generated severe data.

\begin{table}[t]
\centering
\caption{Independent severe-held transfer check (post-campaign, exploratory).
Pooled gain is A5 minus the reference in window-pooled macro-F1 \pp. RMSE
rows apply the frozen coefficients estimated from the original campaign
fitting domain without refitting to the independent severe-held data; the
constant baseline is likewise the frozen campaign-domain constant.}
\label{tab:s5a}
\footnotesize
\setlength{\tabcolsep}{2.4pt}
\begin{tabular}{lrrrrr}
\toprule
Reference & Pooled gain & Sign acc. & RMSE full & RMSE SIR-only & RMSE const.\\
\midrule
IQFormer-insp. & $\SevereTransferGainIqformer$ & \SevereTransferSignAccIqformer
& \SevereTransferRmseIqformerFull & \SevereTransferRmseIqformerSirOnly
& \SevereTransferRmseIqformerConstant\\
MCLDNN & $\SevereTransferGainMcldnn$ & \SevereTransferSignAccMcldnn
& \SevereTransferRmseMcldnnFull & \SevereTransferRmseMcldnnSirOnly
& \SevereTransferRmseMcldnnConstant\\
\bottomrule
\end{tabular}
\end{table}

\section{Independent Severe-Held Validation: Representation Repair by the
I/Q Sidecar}
After the prospectively frozen transfer check revealed a cross-regime transport
boundary on the independently generated severe-held set, a second analysis
plan froze the repair hypothesis, primary contrast, and decision
rule before any sidecar inference on that same held set. The data stay locked
to the same cache; the A5, IQFormer-inspired, and MCLDNN predictions are
reused from the transfer check; and the ten sidecar checkpoints pass a
provenance audit (\SidecarParameters\ parameters each, strict loading of the
frozen state dicts, SHA-256 recorded in
\path{docs/S5_R1_CHECKPOINT_AUDIT.md}). The primary contrast is sidecar
minus A5 under the seed-wise convention, pooled as the equally weighted mean
of the per-seed paired differences, with a seed-level paired bootstrap
interval.

The repair is consistent across seeds and strata. Seed-mean pooled macro-F1
rises from \SevereRepairLevelAFive\%
(A5) to \SevereRepairLevelSidecar\% (sidecar), a difference of
\SevereRepairDiff\ \pp\ ($[\SevereRepairCiLow,\SevereRepairCiHigh]$), with
regime contributions \SevereRepairRegimeOneDiff\ and
\SevereRepairRegimeTwoDiff\ \pp; the difference is positive in
\SevereRepairPositiveSeeds\ of \ProtocolSeedCount\ seeds and in
\SevereRepairPositiveSnrStrata\ of eight SNR strata. A window-level
hierarchical paired audit reproduces the interval. On the same held set the
sidecar also overtakes IQFormer-inspired (\SevereRepairLevelIqformer\%;
descriptive gap-recovery ratio \SevereRepairGapRecoveryIqformer) while
recovering only \SevereRepairGapRecoveryMcldnn\ of the remaining gap to
MCLDNN (\SevereRepairLevelMcldnn\%). The absolute levels remain low for every
evaluated model at this SIR $=-15$~dB operating point, so the verdict is a
consistent relative repair effect rather than deployable classification
accuracy. The repair transfers to the independently generated severe-held
regimes even where the campaign-internal corner advantage and the envelope
transport do not. Per-seed, per-SNR, and per-class tables are stored as CSV
(artifact index below); the sidecar remains exploratory and outside the
sealed family.

\section{Severe-Held Frozen-Inference Attribution Controls}
\label{sec:controls}
A third prospectively frozen analysis plan specified three forward-only
attribution controls on the
same locked cache before any control inference
(\path{docs/S6_SEVERE_HELD_INFERENCE_CONTROLS_PREREG.md}). No model was
retrained, tuned, or re-selected; the controls are post-campaign diagnostic
evidence and cannot change the sealed family or the two earlier severe-held
verdicts.

Control~A forwards the frozen capacity tiers M (\CapacityMParameters\
parameters) and L (\CapacityLParameters\ parameters) on both regimes, with
A5 and the sidecar restricted to the same matched five seeds
(17, 29, 43, 71, 101). Tested spectral width scaling does not reproduce the
repair: M and L reach \SevereControlLevelCapM\% and \SevereControlLevelCapL\%
seed-mean pooled macro-F1 (paired contrasts against A5 of
\SevereControlDiffCapM\ and \SevereControlDiffCapL\ \pp, both intervals
including zero), whereas on the matched subset the sidecar reaches
\SevereControlMatchedLevelSidecar\% against
\SevereControlMatchedLevelAFive\% for A5.

Control~B tests the sidecar's information dependence at inference. B1
replaces the I/Q branch input with zeros of the same shape; B2 applies a
fixed random temporal permutation of the 1024 samples within each window,
identical across the two channels, which preserves the marginal sample
distribution but destroys modulation temporal structure. Both interventions
degrade the sidecar substantially, to \SevereControlLevelIqZero\% and
\SevereControlLevelIqShuffle\%, losses of \SevereControlIqZeroLoss\ and
\SevereControlIqShuffleLoss\ \pp\ relative to the full sidecar.

Control~C removes the residual side-head contribution
$\mathrm{classifier}(\mathrm{fused\ embedding})$ and keeps the co-trained
base spectral logits. The sidecar falls to \SevereControlLevelNoSideHead\%
($\SevereControlDiffNoSideHead$ \pp\ against A5), below A5 itself, so the
added fusion and classifier path is necessary for the repair.

\begin{table}[t]
\centering
\caption{Frozen-inference attribution controls on the severe-held set
(post-campaign diagnostic). Seed-mean pooled macro-F1 in \%. Panels separate
the matched five-seed capacity comparison from the full ten-seed interventions.}
\label{tab:controls}
\footnotesize
\setlength{\tabcolsep}{4pt}
\textbf{(a) Matched five-seed Control A}\\[2pt]
\begin{tabular}{llc}
\toprule
Condition & Seed-mean pooled macro-F1 & Seeds\\
\midrule
A5 & \SevereControlMatchedLevelAFive\% & 5\\
Capacity M & \SevereControlLevelCapM\% & 5\\
Capacity L & \SevereControlLevelCapL\% & 5\\
Sidecar (full) & \SevereControlMatchedLevelSidecar\% & 5\\
\bottomrule
\end{tabular}
\par\vspace{2pt}\parbox{0.96\columnwidth}{\scriptsize All four entries use
exactly the same five algorithm seeds (17, 29, 43, 71, 101); contrasts are
paired within this matched subset.}

\vspace{5pt}
\textbf{(b) Full ten-seed Controls B and C}\\[2pt]
\begin{tabular}{llc}
\toprule
Condition & Seed-mean pooled macro-F1 & Seeds\\
\midrule
Sidecar (full) & \SevereRepairLevelSidecar\% & 10\\
Sidecar, I/Q zeroed (B1) & \SevereControlLevelIqZero\% & 10\\
Sidecar, I/Q time-shuffled (B2) & \SevereControlLevelIqShuffle\% & 10\\
Sidecar, side head removed (C1) & \SevereControlLevelNoSideHead\% & 10\\
\bottomrule
\end{tabular}
\par\vspace{2pt}\parbox{0.96\columnwidth}{\scriptsize Panel (b) uses all ten
frozen sidecar seeds and changes inference only; these results are
information-dependence diagnostics rather than retrained architecture
ablations.}
\end{table}

Under the frozen decision rule---(a) the sidecar exceeds the widest
capacity tier in pooled seed-mean macro-F1; (b) both information losses are at
least $1.0$~\pp; (c) the no-head condition retains at least half of the full
repair contrast---conditions (a) and (b) hold and (c) fails
($\SevereControlDiffNoSideHead$ \pp\ retained). The recorded pattern is
therefore mixed: the repair is best interpreted as a joint
representation-and-fusion intervention rather than an isolated
I/Q-information effect. These are inference interventions, i.e.\ an
information-dependence diagnostic, not an architecture ablation or causal
proof; per-regime levels and paired contrasts are stored as CSV (artifact
index below).

\section{Prospective-Freeze Provenance}
\label{sec:preregp}
Table~\ref{tab:preregp} lists the three frozen analysis-plan documents with their
SHA-256 digests. The transfer check was frozen before any inference on the
new held set; the repair plan was frozen after the transfer result
and before any sidecar inference; the control plan was frozen
after the repair result and before any control inference. Each document
records its primary endpoint and pass/fail rule; the execution logs are kept
alongside the operation logs under \path{docs/}. No commit-level timestamp
service was used; the freeze ordering is established by these documents and
the operation logs rather than asserted beyond that evidence.

\begin{table}[t]
\centering
\caption{Prospectively frozen analysis plans and SHA-256 digests. Freeze order: transfer
check, then repair, then controls; each stage froze before its own inference
on the shared severe-held set.}
\label{tab:preregp}
\footnotesize
\setlength{\tabcolsep}{4pt}
\begin{tabular}{lll}
\toprule
Stage & Document & SHA-256 (first 16 hex)\\
\midrule
Transfer & \path{docs/S5_A_SEVERE_HELD_PREREG.md} &
\texttt{2759558b50a046ea}\\
Repair & \path{docs/S5_R1_PREREG.md} &
\texttt{058cc76dd693b715}\\
Controls & \path{docs/S6_SEVERE_HELD_INFERENCE_CONTROLS_PREREG.md} &
\texttt{be0dc26ccff5b69f}\\
\bottomrule
\end{tabular}
\end{table}

\section{Artifact Index}
\small
Every number in this document is derived from the frozen artifact tree rather
than transcribed from an earlier draft. The relevant entry points are:

\begin{itemize}\setlength{\itemsep}{2pt}
  \item Cell construction and envelope fits:
  \path{analysis_zero_compute/outputs/csv/a14_envelope_cells.csv} and
  \path{a14_v41_envelope_interfered_holdout.csv}.  The V4.4 clean-sentinel
  31-cell sensitivity is recomputed from the same cell table by
  \path{analysis_zero_compute/a16_clean_sentinel_sensitivity.py}, writing
  \path{a16_clean_sentinel_sensitivity.csv} and
  \path{a16_clean_sentinel_sensitivity.json}.
  \item Covariate ladder, selector evaluation, and convergence audit:
  \path{analysis_zero_compute/a17_final_s2_audits.py}, writing
  \path{a17_covariate_ladder.csv}, \path{a17_selector_evaluation.csv}, and
  \path{a17_convergence_audit.csv} from the same frozen cell table and the
  sealed per-seed \path{result.json} training histories.
  \item Leave-$-15$-dB-out test, paired overlap-increment bootstrap, and
  selector confusion matrix:
  \path{analysis_zero_compute/a18_final_s4_audits.py}, writing
  \path{a18_leave15_out.csv}, \path{a18_overlap_increment_bootstrap.csv},
  \path{a18_selector_confusion.csv}, and \path{a18_leave15out_selector.csv}
  from the same frozen cell table and the sealed per-seed predictions.
  \item Sidecar architecture:
  \path{analysis_zero_compute/tier2_gpu/run_tier2_experiment.py}
  (\texttt{LightweightIQSidecarVIMD}); the parameter counts in
  Table~\ref{tab:sidecar} are verified against the sealed checkpoint state
  dict under \path{artifacts/tier2_iq_sidecar_v1/models/}.
  \item Capacity ladder: \path{tier2_capacity_L_v1_vs_S.csv},
  \path{tier2_capacity_L_v1_vs_M.csv},
  \path{tier2_capacity_L_v1_vs_iqformer.csv},
  \path{tier2_capacity_L_v1_vs_mcldnn.csv}, and
  \path{tier2_capacity_M_v2_vs_a5.csv}, all five-seed matched.
  \item Tier-2 interventions: \path{tier2_iq_sidecar_v1_vs_a5.csv},
  \path{tier2_condition_coverage_v1_vs_a0.csv}, and
  \path{tier2_presence_gated_v1_vs_a5.csv}.
  \item Comparator fidelity: \path{artifacts/rml2016_10a_fidelity_v3/}, whose
  \path{run.json} records the dataset SHA-256, the split protocol, the
  reference-release values, and the environment.
  \item The overlap covariate is computed at cache-build time, stored as
  \path{overlap.npy} per split, and SHA-256 checksummed in the sealed cache
  manifest; it is never an input feature of any model.
  \item Independent severe held cache:
  \path{standards/cache_s5_severe_held_v1/}; the manifest records the cache
  digest and the programmatic source-disjointness assertion against the
  campaign source pool.
  \item Severe-held transfer check (Table~\ref{tab:s5a}):
  \path{analysis_zero_compute/a19_s5_severe_held.py}, writing
  \path{analysis_zero_compute/outputs/a19_s5_severe_held.json} together with
  the per-cell CSV files, from the frozen predictions under
  \path{artifacts/s5_severe_held_v1/}.
  \item Severe-held repair check:
  \path{analysis_zero_compute/a20_s5_r1_sidecar_severe_held.py}, writing the
  sidecar predictions under
  \path{artifacts/s5_r1_sidecar_severe_held_v1/} and
  \path{a20_s5_r1_model_levels.csv}, \path{a20_s5_r1_paired_contrasts.csv},
  \path{a20_s5_r1_per_snr.csv}, \path{a20_s5_r1_per_class.csv},
  \path{a20_s5_r1_gap_recovery.csv}, and \path{a20_s5_r1_summary.json}; the
  checkpoint provenance audit is \path{docs/S5_R1_CHECKPOINT_AUDIT.md}.
  \item Severe-held frozen-inference attribution controls
  (Section~\ref{sec:controls}):
  \path{analysis_zero_compute/a21_s6_severe_held_controls.py}, writing the
  control predictions under \path{artifacts/s6_severe_held_controls_v1/} and
  \path{a21_control_levels.csv}, \path{a21_control_matched5_levels.csv},
  \path{a21_control_paired_contrasts.csv}, and
  \path{a21_s6_controls_summary.json}; the frozen analysis plan is
  \path{docs/S6_SEVERE_HELD_INFERENCE_CONTROLS_PREREG.md}.
\end{itemize}

The reproduction package released with the paper contains the resolved macros,
the figure-source CSV files, per-seed aggregates, run configurations, and the
full provenance manifest.

\end{document}
